\documentclass[12pt]{ximera}
\typeout{Start loading xmPreamble.tex}%

% Add here extra macro's that are loaded automatically by all documents of claas 'ximera' or 'xourse' in this repo

%%
%%  Example:
%%
% \newcommand{\R}{\mathbb{R}}
\usepackage{tikz}
\usetikzlibrary{positioning, arrows.meta}
\usepackage{multicol}
\usepackage{enumitem}
\usepackage{caption}
\usepackage{amsmath}
\usepackage{amsthm}
\usepackage{fontenc}

%% ---------------------------------------------------------------
%% Shared worksheet macros.
%%
%% These came from the Summer 2026 worksheet set, where each worksheet carried its
%% own copy. They have to live here instead: a xourse inlines every activity into a
%% single document, so duplicate \newcommand definitions across activities collide,
%% and the xourse gobbles \input so a per-topic macro file never loads.
%%
%% They use \providecommand, NOT \newcommand. ximera.cls loads this file automatically,
%% and every activity also carries an explicit \typeout{Start loading xmPreamble.tex}%

% Add here extra macro's that are loaded automatically by all documents of claas 'ximera' or 'xourse' in this repo

%%
%%  Example:
%%
% \newcommand{\R}{\mathbb{R}}
\usepackage{tikz}
\usetikzlibrary{positioning, arrows.meta}
\usepackage{multicol}
\usepackage{enumitem}
\usepackage{caption}
\usepackage{amsmath}
\usepackage{amsthm}
\usepackage{fontenc}

%% ---------------------------------------------------------------
%% Shared worksheet macros.
%%
%% These came from the Summer 2026 worksheet set, where each worksheet carried its
%% own copy. They have to live here instead: a xourse inlines every activity into a
%% single document, so duplicate \newcommand definitions across activities collide,
%% and the xourse gobbles \input so a per-topic macro file never loads.
%%
%% They use \providecommand, NOT \newcommand. ximera.cls loads this file automatically,
%% and every activity also carries an explicit \typeout{Start loading xmPreamble.tex}%

% Add here extra macro's that are loaded automatically by all documents of claas 'ximera' or 'xourse' in this repo

%%
%%  Example:
%%
% \newcommand{\R}{\mathbb{R}}
\usepackage{tikz}
\usetikzlibrary{positioning, arrows.meta}
\usepackage{multicol}
\usepackage{enumitem}
\usepackage{caption}
\usepackage{amsmath}
\usepackage{amsthm}
\usepackage{fontenc}

%% ---------------------------------------------------------------
%% Shared worksheet macros.
%%
%% These came from the Summer 2026 worksheet set, where each worksheet carried its
%% own copy. They have to live here instead: a xourse inlines every activity into a
%% single document, so duplicate \newcommand definitions across activities collide,
%% and the xourse gobbles \input so a per-topic macro file never loads.
%%
%% They use \providecommand, NOT \newcommand. ximera.cls loads this file automatically,
%% and every activity also carries an explicit \input{../xmPreamble}, so this file is
%% read twice. That was harmless while it held only \usepackage lines (LaTeX skips a
%% package it has already loaded) but a second \newcommand is a hard error.
%%
%%   \blank[width]        fill-in-the-blank rule (default 2.5cm)
%%   \showwork            "show and explain your work" reminder
%%   \showworkline        ... plus which schedule line was used
%%   \showworkyear        ... plus which year's schedule was used
%%   \schedhead{...}      centred bold caption above a tiered-rate table
%%   \samesquares[scale]{left}{right}   two equal-area squares, labelled
%%   \samecubes[scale]{left}{right}     two equal-volume cubes, labelled
%% ---------------------------------------------------------------

\providecommand{\blank}[1][2.5cm]{\underline{\hspace{#1}}}
\providecommand{\showwork}{\textbf{REMEMBER TO SHOW AND EXPLAIN YOUR WORK.}}
\providecommand{\showworkline}{\textbf{REMEMBER TO SHOW AND EXPLAIN YOUR WORK.
  Explanation should include how and why you know which line of the
  schedule was used to calculate this amount.}}
\providecommand{\showworkyear}{\begin{sloppypar}\textbf{REMEMBER TO SHOW AND
  EXPLAIN YOUR WORK. Explanation should include how and why you know which
  line of the year's tax schedule was used to calculate this tax
  amount.}\end{sloppypar}}
\providecommand{\schedhead}[1]{\begin{center}\textbf{#1}\end{center}}

\providecommand{\samesquares}[3][1]{%
\begin{center}
{\footnotesize These squares are the \textbf{same size}.}

\vspace{1.5mm}
\begin{tikzpicture}[scale=#1,every node/.style={scale=#1}]
  \draw (0,0) rectangle (2.4,2.4);
  \node[anchor=north west] at (0.15,2.25) {Area:};
  \node[anchor=east]  at (-0.15,1.2) {#2};
  \node[anchor=north] at (1.2,-0.15) {#2};
  \begin{scope}[xshift=4.6cm]
    \draw (0,0) rectangle (2.4,2.4);
    \node[anchor=north west] at (0.15,2.25) {Area:};
    \node[anchor=east]  at (-0.15,1.2) {#3};
    \node[anchor=north] at (1.2,-0.15) {#3};
  \end{scope}
\end{tikzpicture}
\end{center}}

\providecommand{\samecubes}[3][1]{%
\begin{center}
{\footnotesize These cubes are the \textbf{same size}.}

\vspace{1.5mm}
\begin{tikzpicture}[scale=#1,every node/.style={scale=#1}]
  \foreach \dx in {0,5.2} {
    \begin{scope}[xshift=\dx cm]
      \draw (0,0) rectangle (2.0,2.0);
      \draw (0,2.0) -- (0.65,2.6) -- (2.65,2.6) -- (2.0,2.0);
      \fill[black!12] (2.0,0) -- (2.65,0.6) -- (2.65,2.6) -- (2.0,2.0) -- cycle;
      \draw (2.0,0) -- (2.65,0.6) -- (2.65,2.6) -- (2.0,2.0);
      \node[anchor=north west] at (0.1,1.9) {Volume:};
    \end{scope}
  }
  \node[anchor=east]  at (-0.15,1.0) {#2};
  \node[anchor=north] at (1.0,-0.15) {#2};
  \node[anchor=west]  at (2.25,0.4) {#2};
  \node[anchor=east]  at (5.05,1.0) {#3};
  \node[anchor=north] at (6.2,-0.15) {#3};
  \node[anchor=west]  at (7.45,0.4) {#3};
\end{tikzpicture}
\end{center}}

%% NOTE: do NOT add \usepackage to individual activity .tex files.
%% When a xourse inlines an activity it gobbles \documentclass and \input, but a bare
%% \usepackage still executes -- after \begin{document} -- and errors out.
%% All shared packages and macros belong here.
%%
%% ximera.cls already loads: amsmath, amssymb, amsthm, cancel, enumitem, graphicx,
%% hyperref, listings, pgfplots, tikz, url, xcolor. No need to re-load those.
, so this file is
%% read twice. That was harmless while it held only \usepackage lines (LaTeX skips a
%% package it has already loaded) but a second \newcommand is a hard error.
%%
%%   \blank[width]        fill-in-the-blank rule (default 2.5cm)
%%   \showwork            "show and explain your work" reminder
%%   \showworkline        ... plus which schedule line was used
%%   \showworkyear        ... plus which year's schedule was used
%%   \schedhead{...}      centred bold caption above a tiered-rate table
%%   \samesquares[scale]{left}{right}   two equal-area squares, labelled
%%   \samecubes[scale]{left}{right}     two equal-volume cubes, labelled
%% ---------------------------------------------------------------

\providecommand{\blank}[1][2.5cm]{\underline{\hspace{#1}}}
\providecommand{\showwork}{\textbf{REMEMBER TO SHOW AND EXPLAIN YOUR WORK.}}
\providecommand{\showworkline}{\textbf{REMEMBER TO SHOW AND EXPLAIN YOUR WORK.
  Explanation should include how and why you know which line of the
  schedule was used to calculate this amount.}}
\providecommand{\showworkyear}{\begin{sloppypar}\textbf{REMEMBER TO SHOW AND
  EXPLAIN YOUR WORK. Explanation should include how and why you know which
  line of the year's tax schedule was used to calculate this tax
  amount.}\end{sloppypar}}
\providecommand{\schedhead}[1]{\begin{center}\textbf{#1}\end{center}}

\providecommand{\samesquares}[3][1]{%
\begin{center}
{\footnotesize These squares are the \textbf{same size}.}

\vspace{1.5mm}
\begin{tikzpicture}[scale=#1,every node/.style={scale=#1}]
  \draw (0,0) rectangle (2.4,2.4);
  \node[anchor=north west] at (0.15,2.25) {Area:};
  \node[anchor=east]  at (-0.15,1.2) {#2};
  \node[anchor=north] at (1.2,-0.15) {#2};
  \begin{scope}[xshift=4.6cm]
    \draw (0,0) rectangle (2.4,2.4);
    \node[anchor=north west] at (0.15,2.25) {Area:};
    \node[anchor=east]  at (-0.15,1.2) {#3};
    \node[anchor=north] at (1.2,-0.15) {#3};
  \end{scope}
\end{tikzpicture}
\end{center}}

\providecommand{\samecubes}[3][1]{%
\begin{center}
{\footnotesize These cubes are the \textbf{same size}.}

\vspace{1.5mm}
\begin{tikzpicture}[scale=#1,every node/.style={scale=#1}]
  \foreach \dx in {0,5.2} {
    \begin{scope}[xshift=\dx cm]
      \draw (0,0) rectangle (2.0,2.0);
      \draw (0,2.0) -- (0.65,2.6) -- (2.65,2.6) -- (2.0,2.0);
      \fill[black!12] (2.0,0) -- (2.65,0.6) -- (2.65,2.6) -- (2.0,2.0) -- cycle;
      \draw (2.0,0) -- (2.65,0.6) -- (2.65,2.6) -- (2.0,2.0);
      \node[anchor=north west] at (0.1,1.9) {Volume:};
    \end{scope}
  }
  \node[anchor=east]  at (-0.15,1.0) {#2};
  \node[anchor=north] at (1.0,-0.15) {#2};
  \node[anchor=west]  at (2.25,0.4) {#2};
  \node[anchor=east]  at (5.05,1.0) {#3};
  \node[anchor=north] at (6.2,-0.15) {#3};
  \node[anchor=west]  at (7.45,0.4) {#3};
\end{tikzpicture}
\end{center}}

%% NOTE: do NOT add \usepackage to individual activity .tex files.
%% When a xourse inlines an activity it gobbles \documentclass and \input, but a bare
%% \usepackage still executes -- after \begin{document} -- and errors out.
%% All shared packages and macros belong here.
%%
%% ximera.cls already loads: amsmath, amssymb, amsthm, cancel, enumitem, graphicx,
%% hyperref, listings, pgfplots, tikz, url, xcolor. No need to re-load those.
, so this file is
%% read twice. That was harmless while it held only \usepackage lines (LaTeX skips a
%% package it has already loaded) but a second \newcommand is a hard error.
%%
%%   \blank[width]        fill-in-the-blank rule (default 2.5cm)
%%   \showwork            "show and explain your work" reminder
%%   \showworkline        ... plus which schedule line was used
%%   \showworkyear        ... plus which year's schedule was used
%%   \schedhead{...}      centred bold caption above a tiered-rate table
%%   \samesquares[scale]{left}{right}   two equal-area squares, labelled
%%   \samecubes[scale]{left}{right}     two equal-volume cubes, labelled
%% ---------------------------------------------------------------

\providecommand{\blank}[1][2.5cm]{\underline{\hspace{#1}}}
\providecommand{\showwork}{\textbf{REMEMBER TO SHOW AND EXPLAIN YOUR WORK.}}
\providecommand{\showworkline}{\textbf{REMEMBER TO SHOW AND EXPLAIN YOUR WORK.
  Explanation should include how and why you know which line of the
  schedule was used to calculate this amount.}}
\providecommand{\showworkyear}{\begin{sloppypar}\textbf{REMEMBER TO SHOW AND
  EXPLAIN YOUR WORK. Explanation should include how and why you know which
  line of the year's tax schedule was used to calculate this tax
  amount.}\end{sloppypar}}
\providecommand{\schedhead}[1]{\begin{center}\textbf{#1}\end{center}}

\providecommand{\samesquares}[3][1]{%
\begin{center}
{\footnotesize These squares are the \textbf{same size}.}

\vspace{1.5mm}
\begin{tikzpicture}[scale=#1,every node/.style={scale=#1}]
  \draw (0,0) rectangle (2.4,2.4);
  \node[anchor=north west] at (0.15,2.25) {Area:};
  \node[anchor=east]  at (-0.15,1.2) {#2};
  \node[anchor=north] at (1.2,-0.15) {#2};
  \begin{scope}[xshift=4.6cm]
    \draw (0,0) rectangle (2.4,2.4);
    \node[anchor=north west] at (0.15,2.25) {Area:};
    \node[anchor=east]  at (-0.15,1.2) {#3};
    \node[anchor=north] at (1.2,-0.15) {#3};
  \end{scope}
\end{tikzpicture}
\end{center}}

\providecommand{\samecubes}[3][1]{%
\begin{center}
{\footnotesize These cubes are the \textbf{same size}.}

\vspace{1.5mm}
\begin{tikzpicture}[scale=#1,every node/.style={scale=#1}]
  \foreach \dx in {0,5.2} {
    \begin{scope}[xshift=\dx cm]
      \draw (0,0) rectangle (2.0,2.0);
      \draw (0,2.0) -- (0.65,2.6) -- (2.65,2.6) -- (2.0,2.0);
      \fill[black!12] (2.0,0) -- (2.65,0.6) -- (2.65,2.6) -- (2.0,2.0) -- cycle;
      \draw (2.0,0) -- (2.65,0.6) -- (2.65,2.6) -- (2.0,2.0);
      \node[anchor=north west] at (0.1,1.9) {Volume:};
    \end{scope}
  }
  \node[anchor=east]  at (-0.15,1.0) {#2};
  \node[anchor=north] at (1.0,-0.15) {#2};
  \node[anchor=west]  at (2.25,0.4) {#2};
  \node[anchor=east]  at (5.05,1.0) {#3};
  \node[anchor=north] at (6.2,-0.15) {#3};
  \node[anchor=west]  at (7.45,0.4) {#3};
\end{tikzpicture}
\end{center}}

%% NOTE: do NOT add \usepackage to individual activity .tex files.
%% When a xourse inlines an activity it gobbles \documentclass and \input, but a bare
%% \usepackage still executes -- after \begin{document} -- and errors out.
%% All shared packages and macros belong here.
%%
%% ximera.cls already loads: amsmath, amssymb, amsthm, cancel, enumitem, graphicx,
%% hyperref, listings, pgfplots, tikz, url, xcolor. No need to re-load those.

\title{Rough Estimations: Don't Guess, Round \& Calculate}
\begin{document}
\begin{abstract}
Worked practice in rough estimation: what counts as an estimate rather than a guess,
what a full-credit write-up looks like, five estimation problems to calculate, and then
two passes back over your own work --- deciding whether each estimate is an under- or
overestimate from the rounding you did, and establishing a window of reasonable answers.
\end{abstract}
\maketitle

\section*{Don't Guess: Round \& Calculate}

During today's class session, we will be discussing some scenarios where we would like
to do some rough estimations. It is important to note a few things here.

\begin{enumerate}
\item[(1)] Estimations are NOT guesses. They can sometimes be educated guesses, but
estimations, at least in this course, will always still require some sort of
calculation(s).
\item[(2)] Rough estimations are meant to be estimations which can be calculated
mentally, meaning without calculators, or even pen and paper. You will still be
expected to record/write down your work and thought process for your instructor, but it
should be work that you are able to do mentally.
\item[(3)] Everyone has different skill and comfort levels with their mental
calculations --- you may need to round values a lot to make them something you are
comfortable doing in your head. This is completely fine, but the important thing is
that in your solutions/write-ups you explain what numbers you rounded and why.
\end{enumerate}

For the following problems, do your best to determine a roughly estimated value. Start
off by writing down your assumed/estimated starting values. Be sure to write down all
of your numbers and calculations, indicating anywhere that you rounded. Do not use a
calculator, instead round or estimate values so that you can do all calculations in your
head.

\textbf{NOTE:} Rough estimations are NOT random guesses! They are more like educated
guesses where some ``simple'' calculations are involved in producing the ``guess.''

\subsection*{What a Good Answer Looks Like}

Here are some examples of ways we might expect you to answer rough estimation problems.

For a problem such as: \textit{Roughly estimate how many instant messages you send in a
year (texts, snapchat, whatsapp, etc.)}

\paragraph{Example Solution 1:}
I would estimate that I send about 30 messages a day on average. Some days are more,
and some days are less. At about 30 messages per day, with 365 days in a year, that
would mean an estimate for me would be $30 \times 365$. I am not comfortable
calculating $30 \times 365$ mentally, so I will round 365 to 350.* Then to calculate
$30 \times 350$, I can focus on $3 \times 35$ and then add two 0's to that answer.
$3 \times 35$ is $2 \times 35$ plus another 35, $2 \times 35 = 70$, so
$70 + 35 = 70 + 30 + 5 = 100 + 5 = 105$. Now I add back the two 0's to find that
$30 \times 350 = 10{,}500$. So I would estimate that I send about 10,500 messages per
year.

*Note the rounding choice here could also be 300 or 400 if those are easier for you!

\paragraph{Example Solution 2:}
I would estimate that I send about 500 messages per week on average. At about 500
messages per week, with 52 weeks in a year, that would mean an estimate for me would be
$500 \times 52$. I am not comfortable calculating $500 \times 52$ mentally, so I will
round 52 to 50. Then to calculate $500 \times 50$, I can focus on $5 \times 5$ and then
add three 0's to that answer. $5 \times 5$ is 25, then I add the three 0's to find that
$500 \times 50 = 25{,}000$. So I would estimate that I send about 25,000 messages per
year.

\paragraph{A ``Calculator'' Answer --- would NOT get full credit:}
I would guess that I send about 55 messages per day on average. With 365 days in a
year, that would mean I send $55 \times 365 = 20{,}075$ messages per year. So I would
estimate that I send about 20,000 messages per year.

\paragraph{A ``Math Only'' Answer --- would get at MOST half credit:}
$45 \times 365 = 40 \times 300 = 12{,}000$

12,000 messages per year

\paragraph{An ``Explanation Only'' Answer --- would get at MOST half credit:}
I guessed at how many messages I send in a week and multiplied that by how many weeks
are in a year. I had to round my numbers down to do it mentally, but I got 10,000
messages per year.

\paragraph{A ``Guessing'' Answer --- could get NO credit:}
I would estimate that I send about 15,000 messages per year.

\subsection*{Problem Set 1 --- Calculating}

\begin{problem}
Roughly estimate how much time you spend watching TV shows/movies in a year.

\begin{freeResponse}
\end{freeResponse}

\begin{solution}
I watch maybe one movie a week or a few TV shows per week, so maybe 4 hours per week on
average. There are 52 weeks in a year, so I want to do $4 \times 52$. I am not
comfortable doing that mentally, so I will round 4 to 5 and 52 to 50, making my work
$5 \times 50$. $5 \times 5 = 25$, then include the 0 on the end from the ``50'' to see
that $5 \times 50 = 250$. So I would estimate that I watch about 250 hours of TV a year.
\end{solution}
\end{problem}

\begin{problem}
Roughly estimate how many gallons of gasoline you might use to drive from here to
Chicago, IL (about 320 miles). A car in the US gets an average of 24.4 mpg.

\begin{freeResponse}
\end{freeResponse}

\begin{solution}
To go 320 miles using about 24.4 miles per gallon, and I want to know how many gallons
are used. This means I want to divide 320 miles by 24.4 miles per gallon and this will
give me gallons. I am not comfortable doing $320/24.4$. I will round 24.4 to 25, and
then 320 to 325, so that I figure out $325/25$ for my estimate. 25 goes into 100 four
times, so that means it goes into 300 four $\times$ three times, or 12 times. Then there
is another 25 to get from 300 to 325, meaning 25 goes into 325 a total of 13 times. So
$325/25 = 13$, meaning it would take about 13 gallons to drive from here to Chicago.
\end{solution}
\end{problem}

\begin{problem}
Suppose you were to use an average truck's mpg instead of an average car's mpg. How
would the amount of gallons used change? Explain mathematically why this is the case.

Note: This is not asking for a specific value for how the gallons used would change; it
is asking generally how the value would be different from the previous answer.

\begin{freeResponse}
\end{freeResponse}

\begin{solution}
I would assume a truck would have a lower mpg than a car, meaning we would be dividing
320 miles by a lower number than 24.4 mpg. This would mean we would need more gallons to
go the same amount of miles, as dividing by a smaller number results in a larger number
from the division. So the gallons needed to drive the distance would be higher than for
a car.
\end{solution}
\end{problem}

\begin{problem}
CMH (Columbus airport) averages about 140 flights daily. Roughly estimate how many
flights you think there are in the U.S.\ daily.

\begin{freeResponse}
\end{freeResponse}

\begin{solution}
Suppose each state has about 5 cities that do approximately the same amount of flights
as Columbus. That would be 5 per state $\times$ 50 states $=$ about 250 airports doing
about 140 flights per day. So 250 airports $\times$ 140 flights, that is hard to do
mentally, so I will round 140 to 100, making it $250 \times 100 = $ about 25,000
flights. But there are also some very big airports, like Atlanta, and JFK in New York,
so maybe they do 1,000 flights per day, and maybe there are about 10 of those kinds of
airports in the US, so that would be adding another $10 \times 1{,}000 = 10{,}000$
flights. So I would estimate about 35,000 flights in the US each day.
\end{solution}
\end{problem}

\begin{problem}\label{problem:secondsAlive}
Roughly estimate how many seconds you've been alive.

\begin{freeResponse}
\end{freeResponse}

\begin{solution}
Assuming an age that could be rounded for mental arithmetic to age 20. We would want to
do the ``exact'' calculation of
\[
20 \text{ years} \times 365 \tfrac{\text{days}}{\text{year}}
\times 24 \tfrac{\text{hours}}{\text{day}}
\times 60 \tfrac{\text{min}}{\text{hour}}
\times 60 \tfrac{\text{seconds}}{\text{min}}.
\]

But these are too hard to do mentally, so we'll round various pieces, and maybe even
round some of the results before multiplying by the next value. For example, let's try
rounding 365 days to 360 days, and 24 hours to 25 hours which would give us
$20 \times 360 \times 25 \times 60 \times 60$.

We can multiply numbers together that we know the product for, such as
$20 \times 25 = 500$ and $60 \times 60 = 3600$.

But then we would want the 3600 times the 360, which is too hard to do mentally, so I
will round $3600 \times 360$ to $4{,}000 \times 400 = 1{,}600{,}000$. Then take the 500
from $20 \times 25$ and multiply that by the 1,600,000: $500 \times 1{,}600{,}000$ is
about $500 \times 1{,}500{,}000 = 750{,}000{,}000$.

So: alive about 750,000,000 seconds.
\end{solution}
\end{problem}

\begin{problem}\label{problem:smartphones}
There are about 325 million people in the US. Roughly estimate how many smartphones you
think are sold in the U.S.\ each month.

\begin{freeResponse}
\end{freeResponse}

\begin{solution}
Suppose people tend to get a new phone every 2--3 years on average (some more
frequently, some less frequently). Let's go with every 2 years, since there can be
businesses involved in purchasing phones as well, so that's estimating that about half
the population buys a cell phone each year. Half of 325 is about half of 320 which is
160 million cell phones a year.

We want to divide 160 million by 12 to get monthly, so let's round 12 months to 10.
$160/10 = $ about 16, meaning we estimate about 16 million cell phones are bought each
month.
\end{solution}
\end{problem}

\subsection*{Problem Set 2 --- Reflecting}

In problems like these estimating ones, you might be asked to look at your work and
determine if you think your work is an underestimate or overestimate. In answering such
questions, you should be focusing on your math work to help you decide. Meaning we are
not looking for something like ``well, I think my messages per day is probably less than
what I actually send;'' this would not be a ``math'' reason, but a ``real world''
reason. Instead, you should be looking at the numbers you used in your calculation and
noting whether you rounded UP or DOWN. However, the way you rounded the numbers is not
the only thing you should be noting, you should also include what operation was done
with that rounded value. The problem below will help illustrate why.

\begin{problem}
Multiply 3 times 6, and then multiply 3 times 4. What are the results?

$3 \times 6 = \answer{18}$ \qquad $3 \times 4 = \answer{12}$

If we rounded 6 to 4, we would have been rounding:
\begin{multipleChoice}
\choice{UP}
\choice[correct]{DOWN}
\end{multipleChoice}

By then \emph{multiplying} by this rounded number, we got a:
\begin{multipleChoice}
\choice{LARGER result}
\choice[correct]{SMALLER result}
\end{multipleChoice}

\begin{solution}
$3 \times 6 = 18$ and $3 \times 4 = 12$. Rounding 6 down to 4 and then
\emph{multiplying} by it made the result smaller: $12 < 18$.
\end{solution}
\end{problem}

\begin{problem}
Divide 24 by 6, and then divide 24 by 4. What are the results?

$\dfrac{24}{6} = \answer{4}$ \qquad $\dfrac{24}{4} = \answer{6}$

If we rounded 6 to 4, we would have been rounding:
\begin{multipleChoice}
\choice{UP}
\choice[correct]{DOWN}
\end{multipleChoice}

By then \emph{dividing} by this rounded number, we got a:
\begin{multipleChoice}
\choice[correct]{LARGER result}
\choice{SMALLER result}
\end{multipleChoice}

\begin{solution}
$24/6 = 4$ and $24/4 = 6$. Rounding 6 down to 4 and then \emph{dividing} by it made the
result larger: $6 > 4$.

Compare this with the previous problem. The rounding was identical --- 6 down to 4 ---
but the results moved in opposite directions, because in one case we multiplied by the
rounded value and in the other we divided by it.
\end{solution}
\end{problem}

The operation that was done with a rounded number is just as important as the way the
number was rounded. As such, your explanations should always reference the operation
done as well as how the value was rounded.

Look back at the seconds-alive and smartphones problems from Problem Set 1. Use your
work for those questions to answer the following.

\begin{problem}
For each of those problems, do you think your original rough estimation is an
underestimate or overestimate? Explain why you think this based on your calculations.
Or if you're not sure if it's under or over, explain why you are unsure.

\begin{freeResponse}
\end{freeResponse}

\begin{solution}
For the seconds alive question, I think it is an overestimate, because I rounded
$3600 \times 365$ to $4000 \times 400$, which is multiplying two rounded up numbers that
were rounded up pretty significantly. So although I rounded that result down from
1,600,000 to 1,500,000, I think it was still an overestimate, and that rounding the
result down would not have offset the previous rounding up.

For the smartphones sold each month, mathematically, I rounded 325 million down to 320
million, which could contribute to an underestimate. However, I rounded 12 down to 10
and then divided by that value, which would make a larger result than dividing by 12,
which could contribute to an overestimate. I'm not sure which rounding would have a
larger effect on the resulting value.

If I had to guess, NOT based on the math, I would guess it would be an overestimate
since getting a phone every other year is probably too frequent an estimate, and since I
used the whole 325 million population value, when certainly not all 325 million have a
smartphone (some are elderly, some are infants, some don't want a smartphone, some
aren't allowed, etc.).
\end{solution}
\end{problem}

\begin{problem}
If possible, determine a more exact value for those two problems by using a calculator
and not rounding values to get an exact number, or see if Google has any estimate(s).
Compare these with your original estimations: were you under or over estimating, or
can't tell? If you have values to compare, how different are the two values?

\begin{freeResponse}
\end{freeResponse}

\begin{solution}
For the seconds alive question, an ``exact'' value for a 20 year old:
\[
20 \times 365 \times 24 \times 60 \times 60 = 630{,}720{,}000,
\]
which is less than the 750,000,000 estimate above, meaning I did have an overestimate.

The difference is
\[
750{,}000{,}000 - 630{,}720{,}000 = 119{,}280{,}000 \text{ sec} = 1{,}988{,}000 \text{ min}
= 33{,}133.3\ldots \text{ hours} = 1{,}380.5\ldots \text{ days} = 3.78 \text{ years!}
\]

For the smartphones sold each month, the Statista website estimates about 125 million
smartphone shipments to the US in 2024. Using this as the same as sales, that would be
about $125/12 = $ about 10.42 million per month, which is less than the 16 million per
month estimated in the problem.
\end{solution}
\end{problem}

\subsection*{Problem Set 3 --- Windows of ``Reasonableness''}

Look again at the seconds-alive and smartphones problems from Problem Set 1. Use your
work for those questions to answer the following.

\begin{problem}
Determine a number that would definitely be way too low for even a rough estimate for
each question. Explain why you think this would be unreasonably low. (Still do not just
``guess'' a number, do calculations to get your answer. Still do not use a calculator,
use mental arithmetic.)

\begin{freeResponse}
\end{freeResponse}

\begin{solution}
For the seconds alive question: If we change the $20 \times 365 \times 24 \times 60
\times 60$ to $10 \times 100 \times 10 \times 10 \times 10$ we already get 1,000,000 (1
million) and that's with grossly underestimating the values. Even multiplying that by 10
to get 10,000,000 (10 million) is still a gross underestimate because $2 \times 3 \times
2 \times 6 \times 6$ (the coefficients of each rounded down term of the multiplication)
is definitely more than 10, so multiplying even the 1 million by 10 is still safely an
underestimate.

Definitely too low: 1 million or even 10 million seconds would definitely be too low.

For the smartphones sold each month: If we assume only half of people have smartphones,
about 160 million (from 325 rounded to 320 and dividing by 2), and then assume they only
get a new smartphone every 10 years, which would mean 160 million $/ 10 = 16$ million
phones sold each year. So then divide the 16 million by 12 to get a monthly value, which
could be rounding 16 to 15 million and 12 to 15 months, making it $15/15 = 1$ million
sold each month.

Definitely too low: 1 million sold each month.
\end{solution}
\end{problem}

\begin{problem}
Determine a number that would definitely be way too high for even a rough estimate for
each question. Explain why you think this would be unreasonably high.

\begin{freeResponse}
\end{freeResponse}

\begin{solution}
For the seconds alive question: Thinking similarly to the too low case, too high could
be rounding each value of $20 \times 365 \times 24 \times 60 \times 60$ too high up to
\[
20 \times 400 \times 30 \times 100 \times 100 = 2 \times 4 \times 3 \times 100{,}000{,}000
= 24 \times 100{,}000{,}000 = 2{,}400{,}000{,}000 = 2.4 \text{ billion},
\]
so to get a number this large, we would be rounding each value of the calculation way
too high to be considered reasonable.

Definitely too high: 2.4 billion seconds.

For the smartphones sold each month: If we assume everyone, all 325 million people, gets
a new cell phone each year, then we could divide 325 million by 12. To do so mentally, we
could round 12 down to 10, making an overestimate since we are dividing by a smaller
value. This would give us an estimate of 325 million $/10$ months $=$ about 32.5 million
smartphones sold each month. This seems unreasonably high given the assumption of
everyone getting a new one every year and rounding the months down, making it an estimate
of really more than everyone getting a new phone each year.

Definitely too high: 32.5 million sold each month.
\end{solution}
\end{problem}

\begin{problem}
Collect everyone's estimations for the seconds-alive and smartphones problems on the
board. Are all of the estimates above the ``too low'' and below the ``too high'' values?

\begin{freeResponse}
\end{freeResponse}

\begin{solution}
This will depend on the class's answers! But hopefully all the answers fall within the
``window of reasonableness'' established with the previous questions.

Meaning, hopefully all seconds alive answers are more than 1 or 10 million seconds, and
less than 2.4 billion seconds. Hopefully all monthly smartphone estimates are more than
1 million sold each month and less than 32.5 million sold each month.

If not, it can make for an interesting discussion to figure out why that is not the
case. Was there an error in someone's work? Or is there something reasonable that got a
value outside these bounds?
\end{solution}
\end{problem}

Hopefully, the answer to the final question is ``yes.'' This would mean while everyone
might have different answers based on their individual choices for rounding, that there
is still a window (interval) in which ``reasonable'' estimate answers should fall.

If some of the answers were not within your established window, look more deeply. Were
there some errors in the math done for the values outside the window? Were some
``unreasonable'' rounding choices made? Or perhaps your choices for what would be
unreasonably low or unreasonably high need to be adjusted?

\end{document}
